\section{Residue cancellation, finite index packets and positive denominators}
\label{sec:v41-residue-filters}
The transition formula also identifies when the unmodulated local law
is uniform. It gives a finite-resolution consequence requiring no
arithmetic assumption at all. This consequence concerns a specified
finite packet of indices; it is not substituted for a prescribed single
index in the conditional bridge theorem.

\begin{theorem}[An exact criterion for the unmodulated interval law]
\label{thm:v41-zero-residue-criterion}
Let $K\subset I$ be compact. The following assertions are equivalent.
\begin{enumerate}
\item For every $R\in K$ and every nontrivial occupation resonance,
      $\nu(\eta_Rq_\gamma)=0$.
\item For each $R\in K$, the section phase masses in
      \eqref{eq:v40-section-phase-masses} satisfy $w_{R,l}=c/d_R$.
\item For each fixed positive bounded roof interval $J$, uniformly for
      $R\in K$ and $|Z|\le M$,
\begin{equation}\label{eq:v41-unmodulated-criterion}
 m^2P_{n,m,R}(k,t;J)=c|J|g_{\Omega_R}(Z)+o(1).
\end{equation}
\end{enumerate}
The third assertion is the original $D_R$ Gaussian normalization after
multiplication by $n^2$ as in
\eqref{eq:v40-arithmetic-original-normalization}. It is still an
interval statement, not a pointwise density theorem.
\end{theorem}
\begin{proof}
At fixed $R$, Fourier coefficients of the finite vector
$(w_{R,l})_{l\bmod d_R}$ are the section means of the powers of the
normalized generator. Finite Fourier inversion proves equivalence
of the first two assertions.

Assume the first. Choose the peak charts so that the neighborhood of
the origin is disjoint from every nonzero possible resonance. Such a
choice is uniform by local phase exclusion. On the origin branch,
$a_j=0$, $z_j=1$, $\mu_j=\mu_R$, $H_j=\Omega_R$, and the source
amplitude is $c^2$. Its partition weights sum to one. On every other
chart over $K$, $\kappa_j=0$ implies $A_j(R,a_j)=0$ by the first
assertion. Scalar continuity and compactness, exactly as in
\eqref{eq:v41-damped-curvature-continuity}, give
\[
 \sup_{R\in K}e^{-m\kappa_j(R)}|A_j(R,a_j(R))|\longrightarrow0.
\]
The bounded Gaussian integrals in
\eqref{eq:v41-transition-kernel} show that all these charts contribute
$o(1)$ uniformly, even through changes of resonance type. Theorem~
\ref{thm:v41-uniform-transition} proves the third assertion.

Conversely fix any $R\in K$. Comparison with the fixed-radius
arithmetic law shows $\mathfrak a_R(k,n,m)=1$ on every residue. Every
residue can be tested on a central sequence: take $k=0$, $t=m\bar\tau_R$
and choose $n$ within $d_R$ of $mc$ with $n+h_R^{\rm ar}m$ in the
prescribed congruence class. Then $Z\to0$, and the positive Gaussian
at zero forces that residue factor to equal one. Finite Fourier
inversion of \eqref{eq:v40-finite-residue-sum} makes every nonconstant
squared Fourier coefficient of $(w_{R,l})$ zero. This proves the first
assertion. The argument uses endpoint residues, not triviality of the
whole measurable phase group.
\end{proof}

Let $D_0$ be an integer upper bound for all the orders in
Theorem~\ref{thm:v40-no-roof-resonance}, and fix once and for all
\begin{equation}\label{eq:v41-finite-packet-length}
 L=\operatorname{lcm}\{1,2,\ldots,D_0\},\qquad
 p_L(v)=\frac1L\sum_{j=0}^{L-1}e^{-ijv}.
\end{equation}
This is a finite constant of the mechanical family, not a function of
$m$. No numerical value for the inherited order bound is claimed.

\begin{theorem}[A radius-uniform local law on a fixed finite packet]
\label{thm:v41-finite-packet}
For every fixed positive bounded interval $J$ and $M<\infty$,
\begin{equation}\label{eq:v41-finite-packet}
 \frac{m^2}{L}\sum_{j=0}^{L-1}P_{n+j,m,R}(k,t;J)
        =c|J|g_{\Omega_R}(Z)+o(1)
\end{equation}
uniformly in $R\in I$ and $|Z|\le M$. For all large $m$, the unchanged
union of these $L$ exact-index events has probability at least
$c_{M,J}L m^{-2}$. At least one constituent exact event has probability
at least $c_{M,J}m^{-2}$.
\end{theorem}
\begin{proof}
Sum the exact fixed-count inverse over $n,n+1,\ldots,n+L-1$ and divide
by $L$. This inserts exactly the trigonometric polynomial $p_L(v)$.
On the origin branch its value is one. On a nonzero actual resonance
its value is zero: the occupation coordinate has nontrivial finite
order dividing $L$, by the injectivity of the occupation projection
and the order bound. Thus, on each nonorigin peak chart,
\[
 \sup_R e^{-m\kappa_j(R)}|p_L((a_j(R))_4)|\longrightarrow0.
\]
This is again the compactness argument at the zero set of $\kappa_j$;
it handles nearby nonperipheral branches, not only exact resonances.
The proof of \eqref{eq:v41-test-peak-sum} with the additional smooth
factor $p_L$ consequently leaves just
$c(\int q)g_{\Omega_R}(Z)$, uniformly in the radius. Positive interval
envelopes then give \eqref{eq:v41-finite-packet}, with the band fixed
first. The averaging has been performed on the actual inverse, not
on a fixed-radius formula with an unproved uniform error.

Different genuine return indices cannot share the same collision
time on a single trajectory. These events are therefore disjoint.
Uniform ellipticity and compactness give a positive lower bound for
the right-hand Gaussian on $|Z|\le M$. The two probability assertions
follow. For large $m$, all indices in the packet lie between one and
$m$ because $n/m\to c\in(0,1)$.
\end{proof}

\begin{corollary}[Verified classes for the conditional bridge]
\label{cor:v41-bridge-denominators}
The radius-uniform conditional bridge theorem applies to every central
exact event on a compact set satisfying the equivalent conditions of
Theorem~\ref{thm:v41-zero-residue-criterion}. At each fixed radius it
applies on all positive arithmetic classes. Without either condition,
Theorem~\ref{thm:v41-finite-packet} still guarantees at least one
qualifying exact event in each central packet of the fixed length $L$.
\end{corollary}
\begin{proof}
The respective local laws give the lower mass assumption in
\eqref{eq:v41-denominator-class}. Invoke
Theorem~\ref{thm:v41-actual-return-bridge}. The existence assertion in
a packet does not identify an arbitrary preassigned index with that
index, nor does it replace one conditioning event by another.
\end{proof}

\paragraph{Position in the raw inversion argument.}
The exact-index fixed-interval inverse now has an evaluated expression
uniform through arithmetic transitions, rather than only a uniform
spectral integral. The pinned return path has a conditional Gaussian
bridge whenever the specified microscopic denominator has the stated
lower bound. A fixed finite packet removes all possible occupation
aliases without a diverging resolution scale. The single-index
unmodulated interval law has the exact section-residue criterion above.
None of these assertions sets the remaining section residues to zero
for this family without proof. Nor does the slowly shrinking interval
in Corollary~\ref{cor:v41-shrinking-roof} control density spikes below
that scale. The common pointwise coarea correction and full roof tail
in Theorem~\ref{thm:LLT} remain the requirements for its original raw
density conclusion.
